\chapter{踩坑记录汇编}

本附录汇总开发过程中真实发生的问题。这部分写进方案文档很有说服力——
它们说明团队真的把系统跑通过、压测过，而不只是搭了个架子。

\begin{longtable}{@{}p{2.6cm}p{4.4cm}p{5.4cm}@{}}
\caption{踩坑记录汇编}\label{tab:pits}\\
\toprule
问题 & 表现 & 根因与修法 \\
\midrule
\endfirsthead
\multicolumn{3}{@{}l}{\small 续表 \ref{tab:pits}}\\
\toprule
问题 & 表现 & 根因与修法 \\
\midrule
\endhead
\bottomrule
\endfoot

辩论把数值冲突判成一致 & 两条断言相似度 0.95、出处相同，被判“双方印证”放行 &
数值篡改是文本最相似的一类冲突。判定一致须同时满足文本相似、出处相同、数值集合相同 \\

中文数字绕过数值核验 & “两万小时或五年”改自“一万小时或三年”，全句无阿拉伯数字 &
增加中文数字归一化，不补的话该类幻觉漏掉一半 \\

跨短语二元组误伤真断言 & 跨短语边界拼出的组合被当成特征术语 &
要求两个字的单字文档频率都低于全库比例，高频虚字自然被排除 \\

适配率假性偏低 & 演示只取前四个盲区时掉到 50\%，批量评测却是 92.84\% &
判定漏了知识点难度天花板。\emph{小样本可视化能暴露大样本平均值藏起来的规格问题} \\

追问第一题就触发 & 十六题里七题是追问，覆盖率被吃掉近一半 &
先验低时任何答对都满足条件。加门槛：该知识点必须已有作答证据 \\

学制被当成工龄 & 先验顶到上界，测评开局问偏 &
学制与工龄分开，时长模式加负向前瞻 \\

答案位置严重偏斜 & 四十七道里三十八道在 B，一路选 B 得八成分 &
确定性置换打散，同时重映射画像中的作答下标 \\

判断题全是正命题 & 一路点“正确”就是满分 &
把审核闸反过来用，扰动后必须被判与知识库冲突才采纳 \\

干扰项审核看不见否定 & 手册禁止的做法覆盖率 0.60 被毙，而它正是标准干扰项 &
按分句切开，比较选项更像禁止句还是更像正文 \\

纯数字选项判据失效 & “0.5”切分后只剩两个字符，覆盖率恒为满分 &
数值型与文字型选项分开验证 \\

交叉校验逻辑倒置 & 用相似度当冲突前提，漏掉含义完全不同的两条 &
\emph{越严重的矛盾两条切片越不像}。改为对象相同就比结论 \\

隔离带走已核实的一方 & 注入的错误数据把已核实条目一起挤出知识库 &
冲突双方中已核实的留下；双方状态相同时都留下，叫人来裁 \\

加字段打死整个系统 & 摄入加了新字段，入库后检索与生成全部无法初始化 &
读取端对未知字段容错。知识库由多个工具共写，字段会随工具演进 \\

额度用完炸链路 & 调用异常冒到编排层，演示白屏 &
命题捕获异常返回空值并回退题库 \\

成本显示为零 & 四位小数把单次调用舍进误差 &
改为六位小数 \\

报警代码张冠李戴 & 句式全对、引用格式也对，仅含义安错 &
补术语归属与全库最佳匹配两条规则，这类错误人工审核也容易漏 \\

\end{longtable}

\vspace{6pt}
\noindent
本文档随代码一同交付。代码中的注释记录了更详细的设计理由，
遇到不确定的地方优先看注释——每一处非显然的决定都写明了为什么。
